% =====================================================================================
% AGRADECIMIENTOS, RESUMEN (Revisión 4) Y SUMMARY
% Escribe tu texto FUERA de los recuadros \begin{guia} ... \end{guia}.
% =====================================================================================

\clearpage
\ClearWallPaper
\pagestyle{fancy}

\section*{\centering Agradecimientos}
\addcontentsline{toc}{section}{Agradecimientos}
% >>> ESCRIBE AQUÍ tus agradecimientos (opcional, breve) <<<

A mi familia, por el respaldo constante a lo largo de toda la carrera y por sostener con paciencia las jornadas que exigió este proyecto.

Al Instituto Registral y Catastral del Estado de Tamaulipas, por abrir sus puertas y permitir que este trabajo se desarrollara sobre un problema real de la institución. En particular al Lic. Edgar Francisco Alvarado Tovar, asesor empresarial, por el acompañamiento y la orientación sobre el funcionamiento interno del Instituto.

Al Dr. Said Polanco Martagón, asesor institucional, por la dirección de este trabajo y por las observaciones que le dieron forma, y a los profesores de la Universidad Politécnica de Victoria por la formación recibida a lo largo de la carrera.



\clearpage
\section*{\centering Resumen} % audit:req=resumen
\addcontentsline{toc}{section}{Resumen}

\begin{guia}[Guía: Resumen (Revisión 4)]
\textbf{¿Para qué sirve?} Es la síntesis de todo el trabajo para quien no leerá el documento completo. Se escribe al
final, cuando ya tienes los resultados.

\textbf{Debe responder, en un solo párrafo de 150 a 250 palabras:} ¿cuál es el contexto y el problema? ¿cuál es el
objetivo? ¿qué se hizo (método o solución)? ¿cuáles fueron los resultados principales (con datos)? ¿cuál es la
conclusión o aportación?

\textbf{Reglas:} sin citas, sin figuras, sin siglas sin definir. Cierra con \textbf{4 a 6 palabras clave} (en
minúsculas, separadas por comas).

\textbf{Extensión mínima:} 120 palabras.
\end{guia}

% >>> ESCRIBE AQUÍ el resumen (un párrafo) <<<

El Instituto Registral y Catastral del Estado de Tamaulipas recibe y resguarda un volumen constante de documentación proveniente de sus distintas oficinas. Ese flujo depende del traslado físico de los expedientes, de la entrega presencial y del registro manual de lo recibido, lo que provoca demoras en la atención, riesgo de extravío o duplicidad y, sobre todo, la imposibilidad de conocer con precisión cuánto se está procesando, quién lo procesa y desde qué oficina. El objetivo de este trabajo es integrar un sistema que automatice la solicitud, carga, consulta y control estadístico de documentos sobre un entorno virtualizado, de modo que se reduzcan los tiempos de atención y se garantice la trazabilidad documental. La solución centraliza el proceso en un entorno digital: el usuario registra su solicitud y carga los documentos, que quedan almacenados con sus metadatos y pueden consultarse en línea sin descargarlos, bajo permisos diferenciados según el perfil y la oficina de adscripción del usuario. Cada movimiento queda asentado en una bitácora inmutable que garantiza la trazabilidad del acervo y que constituye, a su vez, la única fuente de datos del módulo de control estadístico, con lo que la información directiva se deriva siempre de los hechos registrados y no de capturas paralelas. El trabajo no se limita al desarrollo del software: comprende también la creación del entorno virtualizado que lo aloja, las políticas de seguridad, respaldo y monitoreo, y su liberación a producción, de manera que el resultado entregado es un sistema operando en un ambiente real y administrado por la propia institución.

\textbf{Palabras clave:} gestión documental, trazabilidad, bitácora de auditoría, administración pública, entorno virtualizado.



\clearpage
\section*{\centering Summary}
\addcontentsline{toc}{section}{Summary}

\begin{guia}[Guía: Summary (Revisión 4)]
Traducción al inglés del Resumen (mismo contenido y estructura), con \textbf{Keywords}. Revisa la gramática; evita
traducción automática literal.
\end{guia}

% >>> WRITE HERE the English summary <<<

The Registry and Cadastre Institute of the State of Tamaulipas receives and safeguards a constant volume of documentation coming from its different offices. This flow depends on the physical transfer of files, in-person delivery and manual recording of what is received, which causes delays, risk of loss or duplication and, above all, makes it impossible to know precisely how much is being processed, who is processing it and from which office. The objective of this work is to integrate a system that automates the request, upload, consultation and statistical control of documents on a virtualized environment, so that service times are reduced and document traceability is guaranteed. The solution centralizes the process in a single digital environment: users register their request and upload the documents, which are stored with their metadata and can be consulted online without downloading them, under permissions differentiated by user profile and office of assignment. Every action is recorded in an immutable audit log that guarantees the traceability of the collection and which is, in turn, the sole data source of the statistical control module, so that management information is always derived from recorded facts rather than from parallel data entry. The work is not limited to software development: it also covers the creation of the virtualized environment that hosts it, the security, backup and monitoring policies, and its release to production, so that the delivered result is a system operating in a real environment managed by the institution itself.

\textbf{Keywords:} document management, traceability, audit log, public administration, virtualized environment.

